\chapter[Spike-train renewal and spectra]{Renewal processes, spike-train spectra, and bin limits}
\label{ch:renewal-spectral}
A spike train may be treated as a random counting measure
$n(t)=\sum_j\delta(t-t_j)$, where the Dirac deltas are
\emph{distributions}, not finite-amplitude voltage spikes.
This representation connects point processes to Fourier spectra
and to continuous-time LFP shot-noise models.

\section{How Bernoulli bins approach a Poisson process}
Partition a fixed duration $T$ into $m=T/\Delta$ bins,
each with independent event probability $p_\Delta=\lambda\Delta$.
Then $N_T\sim\operatorname{Binomial}(m,p_\Delta)$, with
\[
\Prob(N_T=0)=(1-\lambda T/m)^m
\longrightarrow e^{-\lambda T}.
\]
The binomial probability generating function likewise converges
to $\exp[\lambda T(s-1)]$, proving convergence of the
\emph{total count} to a Poisson distribution. Independence
of disjoint bins and consistency across intervals give the
homogeneous Poisson-process limit. A form of Le Cam's inequality
bounds total-variation error in Poisson approximations by a
constant multiple of $\sum_i p_i^2$; hence fixed $T$ and
$p_i=\lambda\Delta$ make this error vanish as $\Delta\to0$.
This approximation is not a guarantee for refractory or bursty
neuronal processes.

\section{Renewal variability}
Suppose interspike intervals $S_1,S_2,\ldots$ are iid,
strictly positive, with mean $m$ and finite variance $v$.
Their coefficient of variation is $CV^2=v/m^2$.
Under standard renewal limit conditions, for long observation
windows $T$ the spike count satisfies
\[
\E[N(T)]\sim T/m,\qquad
\Var[N(T)]\sim (v/m^3)T,\qquad
\frac{\Var N(T)}{\E N(T)}\longrightarrow CV^2.
\]
Thus \emph{long-window} count Fano factor and ISI $CV^2$
are related for a renewal model. Nonrenewal history,
trial mixing, latent gain, short windows and nonstationarity
can break the empirical approximation.

\section{Bartlett spectra and shot-noise LFP}
For a stationary point process, the Bartlett spectrum is the
spectrum of its centered event-counting measure, including
the shot-noise contribution. A homogeneous Poisson train of
rate $\lambda$ has flat continuous spectral density
$S_{nn}(f)=\lambda$ in a two-sided Hz convention
(the DC mean is treated separately).
For a stationary renewal process with inter-event characteristic
function $\varphi(f)=\E[e^{-2\pi ifS}]$, the nonzero-frequency
continuous spectral density takes the form
\[
S_{nn}(f)=\lambda
\frac{1-|\varphi(f)|^2}{|1-\varphi(f)|^2},
\]
under the usual nonlattice and regularity assumptions.
This shows how refractoriness or regular spiking produces
spectral structure without sinusoidal membrane oscillations.

If an idealized LFP component is a linear superposition of
spike-triggered kernels, $y(t)=\int h(t-s)\,dN(s)$,
its centered spectrum satisfies
$S_{yy}(f)=|H(f)|^2S_{nn}(f)$;
independent measurement noise adds its own spectrum.
If both $n$ and $y$ derive from the same events, their cross-spectrum
is proportional to $H(f)S_{nn}(f)$ with sign/conjugation depending
on the cross-spectrum convention. This is an example of
\emph{spike--field coherence without a claim of synaptic causation}.

\section{Hawkes excitation and stability}\label{sec:hawkes}
The classical branching formulation is described by Hawkes
\cite{hawkes1971}.
For a linear Hawkes intensity
$\lambda(t)=\mu+\sum_{t_j<t}h(t-t_j)$ with
$h\ge0$, $\mu>0$ and branching ratio
$\nu=\int_0^\infty h(s)\,ds<1$,
a stationary version exists under the standard construction,
and its mean rate is $\bar\lambda=\mu/(1-\nu)$.
When $\nu\ge1$, this stationary finite-mean construction
fails in the usual nonnegative linear model; the statement
does not apply without alteration to inhibitory or nonlinear
point-process intensities.

\begin{example}[Arsenal interpretation]\sloppy
Compare \path{spike_count_fano.m},
\path{spike_field_coherence.m} and
\path{spike_field_ppc.m} on Poisson,
renewal, and latent-rate simulations. A function operating on
binned counts or spike phases does not itself implement a
continuous-time Bartlett spectral estimator.
\end{example}
% --- expanded edition (October 2026) ---
\section{A central limit theorem for renewal counts}
\begin{theorem}[Renewal CLT]
For a renewal process with interval mean $m$ and variance $v<\infty$,
\[
\frac{N(T)-T/m}{\sqrt{vT/m^3}}\Longrightarrow\Normal(0,1)\qquad(T\to\infty).
\]
\end{theorem}
\begin{proof}[Proof sketch]
Let $S_n=\tau_1+\cdots+\tau_n$. The events $\{N(T)<n\}$ and
$\{S_n>T\}$ coincide. Choose $n=n(T)=\lfloor T/m+x\sqrt{vT/m^3}\rfloor$.
Then
\[
\Prob\Bigl(\frac{N(T)-T/m}{\sqrt{vT/m^3}}<x\Bigr)\approx\Prob(S_n>T)
=\Prob\Bigl(\frac{S_n-nm}{\sqrt{nv}}>\frac{T-nm}{\sqrt{nv}}\Bigr).
\]
Since $T-nm\approx-xm\sqrt{vT/m^3}$ and $nv\approx vT/m$, the right-hand
threshold tends to $-x$, and the ordinary CLT for $S_n$ gives the limit
$\Prob(Z>-x)=\Phi(x)$. Rounding errors in $n$ vanish after scaling.
\end{proof}
The theorem justifies the long-window identity $\Var N(T)\sim vT/m^3$
used for the Fano factor. Its rate of convergence depends on the third
moment of the intervals; for bursty neurons with heavy-tailed intervals
it can be very slow, so the asymptotic Fano factor may not be reached in
windows of a few seconds.

\section{R\'enyi's characterization of the Poisson process}
\begin{theorem}[R\'enyi]
A simple point process on $\R$ is a Poisson process with a diffuse mean
measure $\Lambda$ if and only if $\Prob(N(B)=0)=e^{-\Lambda(B)}$ for every
set $B$ that is a finite union of intervals.
\end{theorem}
The theorem says that the \emph{void probabilities} alone determine a
simple process to be Poisson: one need not check independence of counts
or their full distributions. In practice it motivates a diagnostic:
compute the fraction of windows of length $w$ without spikes and compare
it with $e^{-\hat\lambda w}$ across a range of $w$. For a stationary process with the same rate,
$\Prob(N(w)=0)=1-\bar\lambda w+\Prob(N(w)=2)+2\Prob(N(w)=3)+\cdots$, so
refractoriness and regular firing, which suppress multiple spikes in a
window, \emph{lower} the empty fraction relative to Poisson; bursting,
which packs spikes together, raises it for windows longer than a burst.

\section{Hawkes processes: mean rate and spectrum}
For the linear Hawkes process of Section~\ref{sec:hawkes} with branching
ratio $\nu<1$, take expectations in the intensity:
$\bar\lambda=\mu+\bar\lambda\int_0^\infty h(s)\,ds=\mu+\nu\bar\lambda$,
which gives $\bar\lambda=\mu/(1-\nu)$. Each immigrant event starts a
cluster of expected size $1/(1-\nu)$, matching the branching
interpretation. The Bartlett spectrum is
\[
S_{nn}(f)=\frac{\bar\lambda}{|1-H(f)|^2},\qquad H(f)=\int_0^\infty h(t)e^{-2\pi ift}\,dt.
\]
For an exponential kernel $h(t)=\nu\beta e^{-\beta t}$, $|1-H(f)|^2$ is
smallest at $f=0$, so self-excitation concentrates power at low
frequencies; at $f=0$ the spectrum is $\bar\lambda/(1-\nu)^2$, giving a
long-window Fano factor $1/(1-\nu)^2$. A branching ratio of $0.5$ makes
the Fano factor $4$ even though each event triggers on average only half
an extra event. Self-excitation is thus one mechanism of overdispersion,
alongside latent gain fluctuations (Result 42); the two are hard to
distinguish from counts alone, but differ in their autocorrelograms.

\begin{example}[Gamma renewal spectrum]
For gamma intervals with shape $\kappa$ and rate $\kappa\lambda$ (mean
$1/\lambda$), $\varphi(f)=(1+2\pi if/(\kappa\lambda))^{-\kappa}$. For
$\kappa=1$ (Poisson) the spectrum is flat. For large $\kappa$ the train
is nearly periodic and the spectrum develops peaks at multiples of the
firing rate $\lambda$, with a dip at low frequencies: a regular neuron
transfers little power below its own firing rate. A ``spectral peak at
the firing rate'' in spike--field analysis may simply reflect regular
firing, not an oscillatory input.
\end{example}

\section*{Exercises for this chapter}
\begin{exercise}
Using the renewal CLT, compute an approximate 95\% interval for the
count of a gamma renewal neuron with rate 10\,spikes/s and $\kappa=4$ in
10\,s.
\end{exercise}
\begin{exercise}
Verify R\'enyi's void probability for a homogeneous Poisson process and
show that a dead-time process violates it for short windows.
\end{exercise}
\begin{exercise}
Compute $S_{nn}(f)$ for the exponential Hawkes kernel and find the
frequency at which it falls to half its zero-frequency value.
\end{exercise}
\begin{exercise}
Show that for $\kappa=1$ the gamma renewal spectrum formula gives
$S_{nn}(f)=\lambda$.
\end{exercise}
